\section{Change of expansion for the color sums}\label{L:part2}

The purpose of this section is to represent the finite color sums of
Section~\ref{L:part1} by a convergent Gaussian expansion on a circle.
Proposition~\ref{L:prop:scalar-channel} gives the scalar change of channel
for a fixed list of charges; Section~\ref{L:step5} then recovers the raw
finite sum exactly by a constant coefficient. This representation is the
input to the long-circle estimate in Section~\ref{L:part3}. We first set up
the charge space and justify the formal trace operations, then prove the
scalar identity, and finally establish the stability and convergence
needed to sum over all root insertions.

Use the conjugate convention \(q\mapsto q^{-1}\) in the preceding factors \(P_{ab}\), and write
\[
 b_+=w=\sqrt{k/2}\,e^{i\phi},\qquad b_-=\bar w,\qquad b_0=k .
\]
We use \(\phi=-7\pi/8,-\pi/8\), giving \(r=1,\kappa^2\) respectively. By a raw sum we mean without the multiplier \(\prod K/R^2\); the variables of interest are \(Z=1^N\).

\subsection{Charge space and the formal trace}

Use real vectors \(D_j\), \(j\bmod8\), with one relation \(\sum D_j=0\); take the symmetric bilinear product
\[
 D_0 D_j=(2,-1,1,-2,2,-2,1,-1)_j,\qquad O D_j=D_{j+1}
\]
translated cyclically. It is nondegenerate on this quotient (e.g. cyclic Fourier transform), but not positive. Define \(a\) by \((1-O+O^2)a=D_0\); its row \(aD_j\) is \((1,0,0,-2,0,0,1,0)\). Put \(\Lambda=\sum\mathbb Z D_j,\ S=(1-O)^{-1}\). The color vectors are
\[
 a,\quad a-D_0=(O-O^2)a,\quad a-D_0-D_1=-O^3a ,
\]
all of square \(8/3\). Further put
\[
 \delta=D_7,\quad R_\delta y=y-\delta(\delta y),\quad M_b=O R_\delta,\quad h=-OS\delta=-S D_0,\quad H_b=h^2=9/8 .
\]
The rows \(hD_j=(-1,0,-1,1,-1,1,0,1),\ Sa\cdot D_j=(1,1,1,-1,-1,-1,0,0)\). We have \(S+S^t=1\) (metric transpose), \(S\Lambda\cdot\Lambda\subset\mathbb Z\), \(M_b h=h,\ \delta h=1,\ ah=0\), and
\[
 \det(I-yM_b)=(1-y^6)(1-y^3)/(1-y^2)
\]
on the quotient. These follow by multiplication (one can use \(S=-\sum_{j=0}^7 j O^j/8\)); for the determinant multiply \(\det(I-yO)\) by \(1+\sum_{n\ge1}y^n\delta O^n\delta\) by the rank-one change. In particular the fixed axis of \(M_b\) splits off.

Use parameters, initially \(L>0\),
\[
 p=e^{-L/8},\qquad \mathfrak q=e^{-4\pi^2/L},\qquad s=e^{-2\pi i u/L},\qquad u=x+2\pi i d
\]
with angles \(d\) in an open interval of length 1 about 0. Operators in products act in order of increasing \(d\). The sources (variables \(z=e^{u/8}\)) near 1 have angles close to 0.

The coefficient algebra uses the lattice fields and root-mode relations of vertex-operator constructions \cite[Section~2.1, equations~(2.2)--(2.5)]{Frenkel1985}. We give the formulas and verify their use in this indefinite bilinear space, where convergence is a separate issue. At momentum \(P\), take polynomials freely generated from \(|P\rangle\) by oscillators \(l_{-m}(y)\), \(m>0\), linear in a vector \(y\) in the above space (complexified); let
\([l_n(y),l_m(z)]=n\mathbf1_{m=-n} y z\), positive modes killing the constant, \(l_0(y)=yP\). Degree \(D\) counts oscillator degree (index \(m\) counted \(m\)) plus \(P^2/2\). Traces are coefficient traces on this polynomial space over momentum translates. Use the field
\[
 V_\beta(s)=
 \exp\Big(\sum_{m>0} l_{-m}(\beta)s^m/m\Big)
 \exp\Big(-\sum_{m>0} l_m(\beta)s^{-m}/m\Big)
 T_\beta\, e^{\pi i S\beta\cdot P}s^{\beta P+\beta^2/2},
\]
where \(T_\beta|P\rangle=|P+\beta\rangle\) preserving oscillators. Mode powers count actual changes of \(D\); each specified coefficient has finite output on a polynomial. Momenta at the trace cut range over \({\cal C}=OS(Na)+S\Lambda\). Insert \(\mathfrak q^D e^{-2\pi iD}\widehat O\) at the left, where \(\widehat O\) transforms both momenta and oscillators by \(O\); the source charge shift is \(Na\) up to roots in \(\Lambda\), so the product maps back to this coset. Intermediate momenta always pair integrally with \(\Lambda\).

\subsubsection*{Color shifts and the change of twist}

Zero modes of roots are constant coefficients in \(s\), equivalently contour integrals with \(ds/(2\pi i s)\). For a norm-two root \(b'\) used below, set \(e_{b'}=-[V_{b'}]_0,\ f_{b'}=[V_{-b'}]_0\).

Write \(F_j=f_{D_j}\), \(u_*=\exp(-i\phi)\), and \(g_b=\exp(u_*F_1)\exp(u_*F_0)\).
Our trace \(\mathcal T\) inserts \(g_b V_{Na}(1) g_b^{-1}\) at the fused source, or, for temporarily separated sources, inserts \(g_b V_a(s_i)g_b^{-1}\) at each of them. The commutators below show that the latter field is the sum of the three color fields with coefficients \(1,u_*,u_*^2\); the two addition phases are 1 by the rows above.

Our immediate task is to move the color shifts from the source fields to
the background twist, leaving pure sources. We first verify the root
relations needed for that conjugation, then justify every trace
reordering at finite coefficient support. Analytic summability will be
proved in Section~\ref{L:step6}.

The modes \(e_{b'},f_{b'}\), together with \(l_0(b')\), obey the \(2\)-by-\(2\) traceless-matrix commutation relations (convention \(e\) upper triangular). To check, commuting annihilators to the right in \(V_\beta(s)V_\gamma(s')\) gives \((1-s'/s)^{\beta\gamma}\). Including translations, the ordered product has factor \(e^{\pi i S\beta\cdot\gamma}(s-s')^{\beta\gamma}\) times creation and annihilation exponentials separately combined and the one-site momentum powers on the original momentum (with the combined translation). For a root and a root or source charge the phase commutator supplies \((-1)^{\beta\gamma}\), so moving the root coefficient contour gives ordinary residues. At product \(-1\) the commutator with \([V_\beta]_0\) simply adds \(\beta\) to the other field's charge with coefficient \(e^{\pi i S\beta\cdot\gamma}\), and at nonnegative product it vanishes. Opposite roots have, in \(-V_{b'}(s)V_{-b'}(s')\), the kernel \(s s'/(s-s')^2\) times the expansion starting with
\(1+(s-s')\sum_n l_n(b')s'^{-n-1}\), giving \([e_{b'},V_{-b'}(s')]=\sum_n l_n(b')s'^{-n}\). The reversed-root computation is analogous. These coefficient statements follow also directly from the binomial expansions on the two sides of the pole. Moreover \([l_n(y),V_\beta(s)]=y\beta\,s^n V_\beta(s)\). Thus the triples using the shifted field modes at power \(-n\) and \(l_n(b')\) transform under the adjoint zero-mode action just like at \(n=0\). Both root zero modes are locally nilpotent by degree conservation (shifting repeatedly by the same root eventually raises the momentum degree too far).

For a second expression let \(\widetilde B=\exp(-u_*F_0)\exp(u_*F_1)\exp(u_*F_0)\), \(B=\widehat O\widetilde B\widehat O^{-1}\). Thus \(B\) is an exponential of \(u_* F_2+c u_*^2[V_{-D_1-D_2}]_0\) with commuting summands (the constant \(c\) from the bracket will not be needed). The factor \(\widetilde B\) commutes with \(f_\delta\), and \(B\) with \(e_\delta\), by nonnegative products. Shifting the root indices through \(\widehat O\) gives
\[
 (g_b B)^{-1}\widehat O(g_b B)=\widehat O\exp(-u_* f_\delta)B
   =\widehat M_b X_- B X_+,\qquad X_-=\exp(u_* f_\delta),\quad X_+=\exp(-e_\delta/u_*).
\]
Here \(\widehat M_b=\widehat O W\), where
\[W=\exp(-u_* f_\delta)\exp(e_\delta/u_*)\exp(-u_* f_\delta).\]
This is the root \(SL_2\) matrix with rows \((0,1/u_*),(-u_*,0)\). It takes oscillators to their transforms by \(R_\delta\) (by the adjoint triples), and on constants
\[
 W|P\rangle=(-u_*)^{\delta P}e^{-\pi i(\delta P)S\delta\cdot P}|R_\delta P\rangle .
\]
Indeed for \(n=\delta P\ge0\), \(|P\rangle\) is a highest vector of weight \(n\) by degree, mapped by this matrix to \((-u_*)^n f_\delta^n/n!\) acting on it. This is the elementary degree-\(n\) string action (on \(f^j/j!\) the raising coefficients are \(n-j+1\)). The translation cocycles in the constant component of \(f_\delta^n\) multiply the Laurent constant of \(\prod_j s_j^{1-n}\prod_{i<j}(s_i-s_j)^2\), namely \((-1)^{n(n-1)/2} n!\) by the two alternating determinants, whose sign cancels the cross cocycles. Negative \(n\) uses the analogous lowest vector and \(u_*^n e_\delta^{-n}/(-n)!\).

Conjugate the whole trace by \(g_b B\), preserving the degree insertions. Since \(B,X_+\) commute through the uncolored sources, the result has background \(\widehat M_b\), pure source field(s), and root zero-mode exponentials from \(X_+,B,X_-\). Give them ordered angles \(\epsilon\eta_i\), respectively for
\[
 \delta,\quad -D_2,\quad -D_1-D_2,\quad -\delta,\qquad
 \eta_1<0<\eta_2<\eta_3<\eta_4,\qquad 0<\epsilon\ll1
\]
(with the sources between the first and second). Coefficients and factorials remain those of the indicated exponentials.

\subsubsection*{Finite support at the trace cut}

All trace reorderings here are first at fixed cut degree (degree just before the rightmost operator acts, with the initial state identified with the final one) and fixed powers selected from the source fields, i.e. source mode degrees; in the one-source case only its degree-conserving power enters. They are proper despite the indefinite form. Any cut momentum with a contributing diagonal term, including in cycling the conjugators, must lie in a translate of \(OS\,\operatorname{span}(D_7,D_0,D_1,D_2)\), by the equation after \(\widehat O\) (move left-root shifts inside by \(O^{-1}\), and include the shift by \(\delta\) when using \(W\)). The form on that span after applying \(OS\) is positive: the four-column Gram has diagonal \(9/8\), off-diagonal \(1/8\) except between first and last, \(-7/8\). Thus cut degree bounds the cut momentum in its discrete coset and leaves only finitely many oscillator indices. Successive zero-mode exponentials on each input vector are finite, and fields with prescribed mode degrees have finite outputs. This justifies coefficient composition and cyclicity; in cycling a product past the cut one just matches finite closed matrix paths at the specified degrees (cut degree unchanged since the cycled operators are of degree zero). The root shifts preserve all the relevant lattice translates, as \(\Lambda\subset S\Lambda\). Analytic summability is discussed below.

\subsection{Gaussian change of channel}\label{L:step4}

We specify the scalar evaluation for a fixed list of fields of charges \(\beta_i\) with \(\Delta_{\beta_i}=\beta_i^2/2\) in increasing angular order, either with \(U=O\), or in the transformed trace with \(U=M_b\) (including any root fields in the list). In the latter case it vanishes unless
\(\sum h\beta_i=0\). Set
\[
 b=\sum\beta_i,\quad P_b=OS b,\quad T_U=(1-U^{-1})^{-1}
\]
where \(T_U=1/2\) on a fixed axis. Define \(K_U(Y)\), \(\Re Y>0\), by applying to \(U\) the matrix function with values
\[
 \sum_{m>0,\ m\in\mathbb Z-\nu}\frac{e^{-mY}}m ,
 \qquad y=e^{2\pi i\nu}.
\]
At nonunit-\emph{value} eigenvalues (that is, \(y\ne1\)) take a real lift \(\nu\) modulo 1 continued locally analytically in \(y\), differentiating on Jordan blocks. At 1 simply use positive integers (no Jordan block). Put
\[
 C_U=\lim_{Y\downarrow0}(K_U(Y)+\log Y\,I),\qquad
 d_U(\beta)=e^{-\beta C_U\beta/2}.
\]
The limits (on the positive axis) exist by comparison with \(\sum_{m\in\mathbb Z_{>0}}e^{-mY}/m\): the harmonic coefficient differences are summable and changing a bounded shift in the exponential therein costs \(O(Y(1+|\log Y|))\) locally uniformly in the chosen spectral chart (hence also usable for Jordan derivatives). All the eigenvalues needed here are roots of unity by the characteristic polynomials. The limits are real in the symmetric part, by conjugation and \(U^t=U^{-1}\). For instance if the sequence \(\beta U^j\beta\) has period \(l\) with only corresponding finite-order components, then
\[
 d_U(\beta)=l^{-\Delta_\beta}\prod_{j=1}^{l-1}|1-e^{-2\pi i j/l}|^{(\beta U^j\beta)/2}.
\]

Take real parts \(x_i\) of all representatives in \([-L/2,L/2]\) initially. Sum, for \(M_b\), over \(f\in-\phi/(2\pi)+\mathbb Z\); the \(O\) case has just one term with all fixed-axis contributions below omitted. The phase in units of \(\pi\) is
\[
 \Phi=-\sum_i\Delta_{\beta_i}-\sum_{i<j}\beta_i T_O\beta_j
 -2 f\,\frac{hP_b}{H_b}
 -\sum_{i<j}{\rm sgn}(x_j-x_i)\beta_i T_U\beta_j.
\]
Indices here label angular, NOT \(x\)-order. Include a kinetic factor
\[
 \exp\left[-\sum_{\text{gaps in }x\text{-order}}\frac{\Delta u}{2H_b}
       \left(f+\sum_{\text{before gap}}h\beta_i\right)^2\right],
\]
with cyclic gap lengths in \(u\), value \(f\) of the momentum outside the extremes. For each \(x_i<x_j\) the short-range cross factor has logarithm
\[
 -\beta_i\sum_{n\ge0}\left(K_U(u_j-u_i+nL)
             +K_U(u_i-u_j+(n+1)L)^t\right)\beta_j.
\]
Each field has self multiplier
\[
 (L/(2\pi))^{\Delta_{\beta_i}}d_U(\beta_i)
     \exp\left[-\beta_i\sum_{n\ge1}K_U(nL)\beta_i\right].
\]
In these formulas the transpose of a complex matrix is still the bilinear metric transpose. Finally multiply by
\[
 Z_U=\prod_{l\ge1}\det(I-\mathfrak q^l U)^{-1},
\]
and, for \(M_b\), also \(\sqrt{L/(2\pi H_b)}\). Coefficients from group actions and their factorials are as above, and root integrals now along increasing real coordinate have measure \(du/L\), thus no overall \(L\)-power from root fields. Use limiting boundary values where needed.

\begin{proposition}[Scalar change of channel]\label{L:prop:scalar-channel}
Fix \(L>0\), ordered angular layers with the prescribed common-layer
boundary values, and a finite ordered list of fields
with charges \(\beta_i\) as above. Use the indicated lifts of \(O\) and
\(M_b\) and the cocycle convention of the fields \(V_{\beta_i}\).
For \(U=M_b\), impose the neutrality condition
\(\sum_i h\beta_i=0\); without it the scalar trace contribution is zero.
Let \(\mathcal S_i\) be the displayed self multiplier,
\(\mathcal C_{ij}\) the exponential of the displayed short-range cross
logarithm for \(x_i<x_j\), and \(\mathcal K_f\) the displayed kinetic
factor. The angular trace contribution for this list is
\begin{equation}\label{L:eq:scalar-channel}
 \mathcal T_U(\boldsymbol\beta,\boldsymbol u)
 =Z_U a_U\sum_{f\in\mathcal F_U}
   e^{\pi i\Phi}\mathcal K_f
   \prod_i\mathcal S_i\prod_{x_i<x_j}\mathcal C_{ij},
 \qquad
 a_{M_b}=\sqrt{\frac{L}{2\pi H_b}},\quad a_O=1.
\end{equation}
Here \(\mathcal F_{M_b}=-\phi/(2\pi)+\mathbb Z\).
For \(O\) the sum consists of one term, \(\mathcal K_f=1\), and all
fixed-axis terms in \(\Phi\) are omitted. The identity first holds with
the separated representatives and prescribed boundary values above;
Section~\ref{L:step6} supplies the convergence and continuation required
for the full root expansion.
\end{proposition}

\begin{proof}
The angular cut momentum for \(O\) is \(P_b\); for \(M_b\) it is \(P=P_b+n h\), \(n\in\mathbb Z\), since its equation reads \(P=O(P+b-n\delta)\), \(n=\delta(P+b)\). The constraint is exactly \(hb=0\), and all such \(P\)'s lie in the designated lattice translate. The extra \(n\)-dependence of the position-independent phase before Poisson summation reduces to \(e^{-i n\phi}\). Indeed use the phase
\[
 -P^2+Sb\cdot P-n S\delta\cdot(P+b)
\]
together with \((-u_*)^n\), and \(S\delta=\delta-h\); the difference from \(n=0\) in the display is \(-n^2\). At \(P_b\), including field translation phases gives the first two terms of \(\Phi\). The ordinary momentum factors are
\[
 \exp\left[-\frac{2\pi^2}L P^2-\frac{2\pi i}{L}
       \sum_i u_i\beta_i\Big(P+\sum_{j<i}\beta_j+\beta_i/2\Big)\right].
\]
On the complement of a fixed axis \(P=-T_U b\).

The angular oscillator trace gives \(Z_U\), angular self contractions
\(-\sum_{l>0}\beta_i \mathfrak q^l U(I-\mathfrak q^l U)^{-1}\beta_i/l\), and for \(i<j\), \(Y=u_j-u_i\), the contraction
\[
 -\beta_i\sum_{l\ne0}\frac{e^{2\pi i lY/L}}l
       (I-e^{-4\pi^2 l/L} U^{-1})^{-1}\beta_j.
\]
Indeed at each mode this is just the trace of normal-ordered multiplication and translation exponentials on polynomials, with linear change of variables. Commuting annihilators to the right gives the direct contraction; the normalized trace of the remaining normal product adds contractions from creator vector \(y\) to annihilator \(z\) with matrix \(z\,\mathfrak q^l U(I-\mathfrak q^l U)^{-1} y\) (in unit oscillator normalization). For example diagonalize the linear change: in one direction with eigenvalue including damping equal to \(r\), multiplication by \(e^{a x}\), translation exponential \(e^{b\partial_x}\) in normal order have trace divided by the vacuum geometric series equal to \(e^{ab\,r/(1-r)}\), by summing on \(x^j\). Multiply over directions (using dual annihilators) and continue the resulting matrix inverse to cover a nondiagonal twist, or use its formal series.

For a component \(y=e^{2\pi i\nu}\ne1\), the momentum terms add inside the cross kernel the expression \(2\pi i T_\nu Y/L+4\pi^2 T_\nu(1-T_\nu)/L\), \(T_\nu=1/(1-y^{-1})\). The result on \(0<\Re Y<L\), without the outer minus sign, is the displayed periodic short kernel plus \(\pi i T_\nu\). To check, take Fourier coefficients first for real \(Y\) as boundary values: the short kernel has coefficients, \(k_l=2\pi l/L\),
\[
 \frac1L\sum_{m\in\mathbb Z-\nu}\frac1{m(m+i k_l)}.
\]
Partial fractions (or the elementary cotangent series, by residues of the integer-pole cotangent times the rational fraction) give for \(l\ne0\) exactly \([(1-e^{-2\pi(k_l+i\nu)})^{-1}-T_\nu]/l\), and for \(l=0\) \(4\pi^2 T_\nu(1-T_\nu)/L\). This proves the assertion and gives minus the phase on the opposite interval; it continues analytically with derivatives. On the fixed axis, Poisson summation of the scalar Gaussian gives the square-root prefactor, phase \(-2 f hP_b/H_b\), and exponent \(-L(f-\sum u_i h\beta_i/L)^2/(2H_b)\). The same Fourier comparison, now with \(m\ne0\) integers (subtract the \(m=-\nu\) term then take \(\nu\to0\)), shows the angular oscillator cross kernel minus the short kernel equals
\[
 \frac{Y^2}{2L}-\frac Y2+\frac{\pi i}2-\frac{\pi i}{L}Y+\mathrm{const}
\]
on the positive interval, with sign reversals on the other one (shift by \(L\)). The constant is common and can be counted in self terms by neutrality. Quadratic and imaginary linear corrections thus cancel the Gaussian square correction and the remaining angular-order position term of the momenta. For neutral \(q_i'=h\beta_i\), one has \(\sum_{\rm gaps}\Delta u(\sum_{\rm before}q_i')^2=-\sum_{x_i<x_j}q_i'q_j'(u_j-u_i)\) by summing coefficients across each gap. Thus the remaining terms give precisely the kinetic expression and phase stated above.

These cross checks also determine the self factors: both logs are Gaussian, with any remaining discrepancy a sum of single-charge quadratic forms. Collide a pair \(\beta,-\beta\) from increasing angular and real order (\(Y\to0\)); in the unphased angular normal product the collision factor is \((-2\pi iY/L)^{-\beta^2}\). On the radial side \(-\beta T_U(-\beta)=\Delta_\beta\) in the order phase and the short kernel uses \(-\log Y\,I+C_U\) plus the two self-image sums; this gives exactly that collision using the two proposed self multipliers. Hence there is no discrepancy (one can perform this contraction check for arbitrary continuously varying charges, after stripping the cocycles). All Fourier comparisons continue from \(0<\Im Y<2\pi\) or boundary values from that side. For the real quadratic part, conjugation of the matrix coefficients at a real spectrum of \(\nu\)'s agrees with metric transposition (also for the Jordan derivatives) since \(\bar U=U\), \(U^t=U^{-1}\).

\end{proof}

We will use
\[
 \frac{Z_{M_b}\sqrt{L/(2\pi H_b)}}{Z_O}
 = e^{7L/192}\frac{P_2 P_8}{P_6 P_3 P_1},\qquad
 P_j=\prod_{n\ge1}(1-e^{-nL/j}).
\]
Indeed if \(F(t)=\sum\log(1-e^{-nt})\), then \(F(t)=F(4\pi^2/t)+t/24-\pi^2/(6t)+\frac12\log(2\pi/t)\). For one quick proof apply ordinary Poisson summation to \(|x|/(e^{t|x|}-1)\). Its Fourier coefficients at integers are the samples of the transform \(t^{-2}(y^{-2}-\pi^2/\sinh^2(\pi y))\), \(y=2\pi l/t\), obtained by summing the exponential series and the elementary cotangent partial fractions. They are absolutely summable. This gives \(1/t+2F'(t)=\pi^2/(3t^2)+1/12-8\pi^2 F'(4\pi^2/t)/t^2\); integrate using \(t=2\pi\). The formula for the determinant ratio follows using the characteristic polynomials above and \(\det(I-yO)=(1-y^8)/(1-y)\).

\subsection{Recovery of the finite sums}\label{L:step5}

The purpose of this normalization is exact coefficient recovery. We first
fix the source multiplicity $N$ and the real period, and fuse the sources;
Section~\ref{L:part3} will separately control the fused charge $Na$ as
$N$ grows. Take \([p^0]\) at zero (Taylor, or Cauchy average) of the
following after confluence:
\[
 \mathscr P_N= w^N
 \frac{\mathcal T}{Z_O\,e^{-\pi i\Delta_{Na}}d_O(Na)(L/(2\pi))^{\Delta_{Na}}}.
\]
This gives the raw color sum. To check, initially separate the \(N\) sources near 1 (distinct angular and real coordinates), and divide \(\mathcal T/Z_O\) instead by the leading uncolored radial prescription there at \(O\), i.e. without the finite-period images. The \(d_O\) ratios for the middle color and last color relative to the first are \(\sqrt{2k}\) and 1 by the self formula. Pair phase differences from uncolored are even integral (\(S\) pairings change by integers). For inner/outer fields \(y,z'\) in radial order, ratio of inner to outer source variable \(\zeta\), the direct short kernel exponentiates to
\[
 B_{z'y}(\zeta)=\prod_{j=0}^7(1-q^{-j}\zeta)^{z'O^j y}.
\]
Its ratio for colors versus \(a,a\) is \(P_{\alpha_{\rm inner},\alpha_{\rm outer}}(1/\zeta)\) in conjugate convention by substitution. Thus the zero-image sum with the displayed multiplier has exactly the wanted color weights (including \(w,u_*\) and the self ratios). Positive images add analytic factors starting with 1 near \(p=0\), using arguments multiplied by \(p^n,\ n\ge1\).

For fixed real period the fusion limit of this separated ratio gives the displayed ratio. Indeed before moving the conjugation into colors (or on the transformed side with pure \(V_a\)'s, root contours away), the leading local angular contraction of the \(N\) uncolored fields combines them to \(V_{Na}\); its scalar, from the common powers, translation phases and direct pair factors, is the same as in the pure uncolored angular trace. All other contractions are regular there and merge multiplicatively. In passing from that pure angular trace to the leading radial denominator we only strip finite image factors, smooth and merging the same way. This argument can first use small \(\mathfrak q\) with the absolute bounds below: in the transformed angular contractions at separated sources, factor off the common direct uncolored collision powers before estimating. After this cancellation the remaining Taylor/image factors are regular with the same uniform bounds through coincidence (root layers strictly away and fixed); opening/closing the source separations only changes lower-order terms of the angular-momentum quadratic in root counts and the axis index. Then both the fused transformed radial expression and the separation coefficient of the color expression used next continue along real positive periods. Equivalently one can fuse in the continued radial expressions with root contours still away. To commute fusion with \([p^0]\), let the separated sources approach 1 along generic fixed log-slopes times a small scalar. The normalized \(O\)-color expression is meromorphic in that scalar with bounded pole orders and holomorphic image series in a common disk of \(p\). Its fusion value is its constant Laurent coefficient (negative coefficients vanish by the fusion limit on real periods). Thus the recovered term is exactly the previously established homogeneous rational color sum, not a relative approximation in \(p\).

\subsection{Stability and analytic justification}\label{L:step6}

We now justify summing the transformed scalar formula over all root
insertions. There are two distinct tasks. At small $\mathfrak q$ we compare
convergent angular expansions with the coefficient identities already
proved. We then use a bound on the radial interactions to continue that
identity to every positive period and to the bounded complex shifts of
large periods needed in Section~\ref{L:part3}. Source multiplicity is fixed
throughout this continuation and through the confluence argument of
Section~\ref{L:step5}; uniform estimates for the fused charge $Na$ will be
proved separately in the long-circle argument.

\subsubsection*{A positive covariance for the four root layers}

The full boson form is indefinite. The positivity we need concerns only
the real part of the short-range covariance on the four root layers,
excluding the order phases and kinetic factors. Write their charges as
\[
 \gamma_i=\sigma_i r_i,\qquad
 r=(\delta,D_2,D_1+D_2,\delta),\qquad
 \sigma=(1,-1,-1,-1),\qquad v_i=hr_i.
\]
Fix $\eta_1<0<\eta_2<\eta_3<\eta_4$, and put the roots on layers
$d_i=\epsilon\eta_i$. We will choose one sufficiently small
$\epsilon>0$ and keep it fixed. All constants below may depend on these
layers; no estimate uniform as $\epsilon\downarrow0$ is required.

For a real period $L$, set $t'=4\pi^2|l|/L$ and $y=e^{-t'}$. The Fourier
matrix of the unsigned covariance at a nonzero frequency is $2\pi^2/L$
times $A_\epsilon(t')/t'$, where
\[
\begin{split}
 A_\epsilon(t')&=P-\frac{2vv^{\mathsf T}}{t'H_b},\\
 P_{ij}&=r_i\left(e^{-dt'}(I-yM_b^{-1})^{-1}
       +e^{dt'}yM_b(I-yM_b)^{-1}\right)r_j,
 \qquad d=\epsilon(\eta_j-\eta_i),\quad i\le j,
\end{split}
\]
and $P$ is made symmetric. Here $\mathsf T$ is ordinary transpose on the
four-component column array. The two angular oscillator terms give $P$:
their coefficients are real, so their real parts contribute cosines.
The fixed-axis Fourier comparison in the scalar-channel proof subtracts
$2vv^{\mathsf T}/(t'H_b)$. At frequency zero we use the continuous limit
of $A_\epsilon(t')/t'$. This is also the correct zero Fourier coefficient:
the line kernel has an integrable local logarithmic singularity and
exponential tails, and the nonzero sample formula holds for arbitrary
real periods, hence identifies its continuous line Fourier transform.

We claim that, after choosing $\epsilon$ sufficiently small, there is
$c_\epsilon>0$ such that
\begin{equation}\label{L:eq:root-fourier-coercivity}
 A_\epsilon(t')\ge c_\epsilon\min(t',1)I
 \qquad(t'>0).
\end{equation}
The verification has three frequency regimes. First set $\epsilon=0$
and retain only the three distinct rows. Put
\[
 b_y=y^2+y+1,\qquad a_y=y^2-y+1,\qquad
 e_y=3y^4+3y^3+4y^2+3y+3.
\]
Direct substitution gives the three leading principal minors of $P$ and
the contraction needed for the rank-one subtraction:
\[
\begin{aligned}
 &\frac{2(y+1)(y^2+1)(y^4+1)}{(1-y)a_yb_y^2},
 \quad\frac{y(y+1)^2(4y^4+7y^3+10y^2+7y+4)}{b_y^4},
 \quad\frac{2y(1-y)(y+1)^3e_y}{a_yb_y^4},\\
 &z_y=\frac{2v^{\mathsf T}P^{-1}v}{H_b}
       =\frac{32(1-y)b_y^2}{9(y+1)e_y}< -\log y.
\end{aligned}
\]
Indeed the derivative of $-\log y-z_y$, multiplied by
$9y(y+1)^2e_y^2$, is minus $(1-y)^2$ times the palindromic polynomial of
degree eight with first five coefficients $81,390,1145,2042,2476$;
the difference vanishes at $y=1$. Thus the three-row subtracted matrix
is strictly positive for $t'>0$. Writing $A_0^{(3)}$ for this three-row principal block, its removable
small-frequency limit is
\[
 \lim_{t'\downarrow0}\frac{A_0^{(3)}(t')}{t'}
 =\frac1{27}\begin{pmatrix}
 44&-26&10\\-26&32&8\\10&8&44
 \end{pmatrix}>0.
\]
The leading principal minors of this limit are $44/27$, $244/243$ and
$272/243$. On restoring the fourth row, the sole additional zero direction
at $\epsilon=0$ is the duplicate difference. Since
$\partial_dP_{14}|_{d=0}=-2t'$, separating its two layers lifts this
direction positively to first order in $\epsilon$. After division by
$t'$ these expansions are removable at zero. They establish
\eqref{L:eq:root-fourier-coercivity} on each fixed bounded frequency
interval, including zero, for sufficiently small fixed $\epsilon$.

A separate argument supplies uniformity at high frequency. At
$\epsilon=0$ the matrix approaches the positive semidefinite Gram
matrix $C=(r_ir_j)_{ij}$; on its first three rows the upper triangle is
$2,-2,-1;2,1;2$. The vector $v$ is orthogonal to $\ker C$. In the scale
$s'=\epsilon t'$, the layer perturbation has derivative
$(-C_{ij}|\eta_j-\eta_i|)$. Represent $C$ by Euclidean vectors $c_i$.
For $z\in\ker C$,
\[
 -\sum_{i,j}z_iz_jC_{ij}|\eta_j-\eta_i|
 =2\sum_{j=1}^3(\eta_{j+1}-\eta_j)
       \left|\sum_{i\le j}z_ic_i\right|^2>0
 \quad(z\ne0).
\]
To make this lift uniform in the crossover, fix a large $T$ independently
of $\epsilon$. For $t'\ge T$, the positive semidefinite matrix
$A_0(t')$ has two eigenvalues bounded below by a fixed $a>0$; its
complementary two-dimensional eigenspace $E_{t'}$ tends to $\ker C$.
In the variable $s'$ we have, uniformly for $t'\ge T$,
\[
 A_\epsilon(t')=A_0(t')+s'H(t')+O(s'^2),
 \qquad H(t')\longrightarrow(-C_{ij}|\eta_j-\eta_i|).
\]
Consequently $H(t')$ is bounded below by a fixed $b>0$ on $E_{t'}$
when $T$ is large. Decompose a vector into $p\in E_{t'}^\perp$ and
$q\in E_{t'}$. The unperturbed quadratic form is at least
$a|p|^2$ and is nonnegative on $q$. Absorb the possible
$O(s')|p|^2$ term and the mixed $s'H(t')$ term in half of
$a|p|^2$, at cost $O(s'^2)|q|^2$; for sufficiently small
fixed $s_0>0$, this cost and the Taylor remainder are absorbed in
$b s'|q|^2$ whenever $0<s'\le s_0$. This gives a lower bound
$c(|p|^2+s'|q|^2)$ uniformly in this regime. In particular, no comparison
between $e^{-t'}$ and $s'$ is needed: the full positive semidefinite
$A_0(t')$ was retained.

For $s'\ge s_0$, the entrywise product
$C_{ij}e^{-s'|\eta_j-\eta_i|}$ is strictly positive definite, by the
positive Fourier transform of the exponential covariance, and tends to
$2I$ at infinity. Its least eigenvalue is uniformly positive on this
range. The rank-one subtraction is $O(1/t')$, and the remaining errors
from $y$ decay exponentially, uniformly when
$\epsilon|\eta_j-\eta_i|<1/2$. Since $t'\ge s_0/\epsilon$, both
errors are uniformly small after decreasing $\epsilon$. Thus the order
of choice is $T$, then $s_0$, then a sufficiently small $\epsilon$.
These high-frequency
bounds overlap the bounded-frequency estimate and prove
\eqref{L:eq:root-fourier-coercivity} for all $t'>0$ with one fixed
$\epsilon$.

\subsubsection*{From the covariance to an unsigned root-count bound}

Partition the real circle into successive bins of lengths between
$1/2$ and $2$. Let $K_{I,i}$ count roots of species $i$ in bin $I$, and
put $K_I=\sum_iK_{I,i}$ and $K=\sum_IK_I$. We claim that the product of
root short-range factors, including finite self images and the constants $d_U(\gamma_i)$,
satisfies
\begin{equation}\label{L:eq:root-bin-stability}
 \left|\text{root short-range product}\right|
 \le \exp\left(C_0K-c_0\sum_IK_I^2\right).
\end{equation}
The constants are uniform for all sufficiently large real periods.
The product here excludes the factors $L/(2\pi)$ from root self
multipliers and the coordinate measures; those pair to give $du/(2\pi)$
per root in the integrated formula.

To obtain a bounded kernel from the logarithmic covariance, split the
diagonal scalar symbol $2/(t'+a_*)I$ from $A_\epsilon(t')/t'$, with
$a_*$ sufficiently large. At infinity the latter symbol is
$2I/t'+O(t'^{-2})$, and away from infinity use
\eqref{L:eq:root-fourier-coercivity}. The remainder is therefore an
integrable symbol with a bounded continuous spatial kernel, and it
still dominates $c/(1+t'^2)I$. The removed scalar kernel has nonnegative
spatial values away from the diagonal: use
$2/(t'+a_*)=2\int_0^\infty e^{-s(t'+a_*)}\,ds$. It consequently causes no growth in
pairs of distinct points of the same species. We may drop those factors
when proving an upper bound.

The charge signs do not change the lower comparison. If
$\rho_i=\sum_{\text{species }i}\delta_x$ and
$D_\sigma=\operatorname{diag}(1,-1,-1,-1)$, the covariance is conjugated by $D_\sigma$,
while the scalar diagonal comparison is unchanged. Its quadratic energy
is bounded below by
\[
 c\sum_i\iint G_L(x-y)\,d\rho_i(x)d\rho_i(y)
 \ge c'\sum_{I,i}K_{I,i}^2
 \ge\frac{c'}4\sum_IK_I^2,
\]
where $G_L$ is the positive periodized exponential kernel corresponding
to $1/(1+t'^2)$. The first inequality uses its uniformly positive value
on distances within one bin. Insert and then remove the diagonal terms
of the bounded remainder at a cost $O(K)$. The singular scalar diagonal
was never inserted. Finite self images and the constants $d_U$ also cost
$O(K)$. This proves \eqref{L:eq:root-bin-stability}, including for the
signed charges. In particular coincident same-species logarithmic
singularities cause no absolute-value problem. On compact subintervals
of $L>0$, use bins scaled to the period if necessary; the same bound
holds with locally adjusted constants.

\subsubsection*{Equality at small $\mathfrak q$ and continuation in the period}

We first compare the two analytic expressions at small
$\mathfrak q>0$, retaining the distinct angular layers and, initially,
separated sources. At their radii $|s_i|=\mathfrak q^{-d_i}$, contract
the free oscillators for each fixed list of fields. For $K$ roots and
axis index $n$, the angular zero momenta give a power of $\mathfrak q$
equal to the time-average of the momentum norms divided by two. Its
homogeneous quadratic part in the counts and $n$, subject to $hb=0$,
is at $\epsilon=0$ the cut norm divided by two on the positive space
$OS\operatorname{span}(D_7,D_0,D_1,D_2)$ used above. The norms before
and after the collapsed root product agree by twist invariance. Its only
kernel in the root-count cone consists of equal numbers of $\delta$
and $-\delta$, with $n=0$: a vanishing cut momentum forces the root sum
to be $n\delta$, and neutrality then forces that sum to vanish.
Opening the interval between these two layers lifts the remaining
direction positively to first order, since $\delta^2=2$. With our fixed
sufficiently small $\epsilon$, the exponent is therefore
\[
 \ge c(K^2+n^2)-O(1),
\]
where constants may depend on the fixed external fields.

For each list and momentum, the formal oscillator coefficients can now
be combined by the contraction formula. Their absolute Laurent-series
norms on the layer contours cost at most $\exp(O(K^2+1))$. Direct ratios
are strictly inside their annuli except for equal-layer contacts of
identical roots, whose factors are integer-power polynomials. Image
series converge geometrically even with the Jordan factors, and $Z_U$
has a convergent series as well. These bounds also justify extracting
the specified root coefficients on the contours. The quadratic
$\mathfrak q$ decay dominates these costs when $\mathfrak q$ is small.
The endpoint $O$-formula converges with separated sources. Together with
the earlier finite coefficient-support argument, this proves equality
of the analytic expressions there by grouping coefficients.

For continuation along the positive $L$ axis, use
\eqref{L:eq:root-bin-stability}. At fixed $L$ and fixed source
multiplicity, separated external fields cost only $\exp(O(K))$ in the
root bound, and the momentum sum has Gaussian tails for
$|f|\gg K+1$. Factorials from the root exponentials then give absolute
summability. Since root layers remain separated from the sources,
source confluence at fixed $L$ is harmless after removing the common
uncolored collision factors as in Section~\ref{L:step5}. This continues the equality of the two normalized trace expressions
to every positive real period. The raw finite sum is then recovered by
$[p^0]$ as in Section~\ref{L:step5}; neither this continuation nor the
fixed-multiplicity fusion asserts uniformity in the source multiplicity.

We finally establish the complex-period bounds used by Cauchy averaging.
After fixing the layers and their stability margin, choose a common bend
\[
 u=x+i\alpha(x)+2\pi i d,\qquad
 \alpha(x+\Re L)=\alpha(x)+\Im L,\qquad\alpha(0)=0,
\]
whose Lipschitz constant is sufficiently small relative to $c_0$.
For any prescribed bounded range of $\Im L$, the bend may be confined
to a sufficiently large fixed width at the seam when $\Re L$ is large.
For small real separation $y'$, it changes the imaginary displacement
by $O(\|\alpha'\|_\infty|y'|)$, whereas the kernel derivative is
$O(1/|y'|)$. At large distance the derivative decays exponentially,
including its polynomial Jordan factors. Summing the changes over bins
therefore costs at most a small multiple of $K+\sum_IK_I^2$ in the real
logarithm, which is absorbed in \eqref{L:eq:root-bin-stability}.
The modulus of the kinetic factor still uses the real gap lengths.

For holomorphy, fix $L_0>0$ and parameterize all contours by
\[
 u=x+\theta(x)(L-L_0)+2\pi i d,
 \qquad -L_0/2\le x\le L_0/2,
\]
where $\theta$ is real and piecewise smooth,
$\theta(x+L_0)=\theta(x)+1$, $\theta(0)=0$, and
$\theta(\pm L_0/2)=\pm1/2$. For an arbitrary fixed $L_0$, restrict to a sufficiently small complex
neighborhood so that the derivative times the imaginary shift is small.
For large $L_0$, a large fixed seam width instead accommodates any
prescribed bounded imaginary shift range with the same small slope.
For real $L$ near $L_0$ this monotonically reparameterizes straight
contours. For the nearby complex periods under consideration, signs of
real-coordinate order are still determined by $x$. On each fixed sorted
chamber the integrands, with their continued branches, are holomorphic
almost everywhere in $L$ and are locally dominated by the summable
bounds just proved, comparing with the curves at the real-part period.
Integration thus gives the claimed continuation, moving external
positions with the contours when necessary. There is no need to refold
by the order of $\Im(u/L)$ after bending. Source coincidences use exactly
the regular fixed-multiplicity fusion limits of Section~\ref{L:step5}.
